\chapter*{研究现状图片预览}
\addcontentsline{toc}{chapter}{研究现状图片预览（临时）}

% ============================================================
% 统一样式定义
% ============================================================
\tikzset{
    % 基础模块样式
    module/.style={
        draw, rounded corners=2pt, minimum height=0.7cm,
        font=\small, align=center, thick,
        fill=#1!12, draw=#1!70!black
    },
    module/.default=blue,
    % 箭头样式
    arrow/.style={-{Stealth[length=5pt]}, thick},
    arrow2/.style={-{Stealth[length=4pt]}, semithick},
    dasharrow/.style={-{Stealth[length=5pt]}, thick, dashed},
    % 大括号标注
    brace/.style={decorate, decoration={brace, amplitude=5pt, mirror, raise=2pt}},
    braceup/.style={decorate, decoration={brace, amplitude=5pt, raise=2pt}},
    % 标签
    lbl/.style={font=\footnotesize, align=center},
    slbl/.style={font=\scriptsize, align=center},
    % 虚线框
    dbox/.style={draw, dashed, rounded corners=3pt, inner sep=6pt, #1},
    dbox/.default={gray!60},
}

% ============================================================
% 图 2-1：建模范式对比图
% ============================================================
\begin{figure}[htbp]
\centering
\begin{tikzpicture}[node distance=0.5cm and 0.5cm]
    % 固定左右输入输出的 x 坐标，确保三行对齐
    % 总宽控制在 ~12cm 以内（A4 textwidth 安全范围）
    \def\xleft{-0.4}    % 输入语音 x
    \def\xright{11.4}   % 输出语音 x
    \def\xcenter{5.5}   % 图表中轴线

    % (a) 三框：ASR 中心=2.0, MT 中心=5.5, TTS 中心=9.0
    % ASR: 宽2.0, 左=1.0, 右=3.0
    % MT:  宽2.4, 左=4.3, 右=6.7 (双行，高度自然撑大)
    % TTS: 宽2.0, 左=8.0, 右=10.0
    % 箭头间距: 3.0→4.3=1.3cm, 6.7→8.0=1.3cm

    % === (a) 三阶段级联系统 ===
    \def\ya{0}
    \node[module=green, minimum width=2.0cm, minimum height=0.9cm] (asr) at (2.0, \ya) {语音识别};
    \node[module=orange, minimum width=2.4cm, minimum height=0.9cm] (mt) at (\xcenter, \ya) {机器翻译/\\[-2pt]大语言模型};
    \node[module=red, minimum width=2.0cm, minimum height=0.9cm] (tts) at (9.0, \ya) {语音合成};
    \node[font=\small] (sin1) at (\xleft, \ya) {输入语音};
    \node[font=\small] (sout1) at (\xright, \ya) {输出语音};
    \draw[arrow] (sin1) -- (asr);
    \draw[arrow] (asr) -- node[above, yshift=1pt, font=\scriptsize]{转写文本} (mt);
    \draw[arrow] (mt) -- node[above, yshift=1pt, font=\scriptsize]{目标文本} (tts);
    \draw[arrow] (tts) -- (sout1);
    \node[font=\small\bfseries] at (\xcenter, \ya-1.0) {(a) 三阶段级联系统};

    % === (b) 两阶段级联系统 ===
    % S2T 框：左边缘对齐 ASR 左=1.0，右边缘对齐 MT 右=6.7
    % 宽度=5.7cm, 中心=(1.0+6.7)/2=3.85
    % TTS 框：与 (a) 的 TTS 对齐，中心=9.0, 宽=2.0
    \def\yb{-2.5}
    \node[module=teal, minimum width=5.7cm, minimum height=0.9cm] (s2t) at (3.85, \yb) {语音到文本模型};
    \node[module=red, minimum width=2.0cm, minimum height=0.9cm] (tts2) at (9.0, \yb) {语音合成};
    \node[font=\small] (sin2) at (\xleft, \yb) {输入语音};
    \node[font=\small] (sout2) at (\xright, \yb) {输出语音};
    \draw[arrow] (sin2) -- (s2t);
    \draw[arrow] (s2t) -- node[above, yshift=1pt, font=\scriptsize]{目标文本} (tts2);
    \draw[arrow] (tts2) -- (sout2);
    \node[font=\small\bfseries] at (\xcenter, \yb-0.8) {(b) 两阶段级联系统};

    % === (c) 端到端模型 ===
    % E2E 框：左边缘对齐 ASR 左=1.0，右边缘对齐 TTS 右=10.0
    % 宽度=9.0cm, 中心=(1.0+10.0)/2=5.5
    \def\yc{-4.5}
    \node[module=violet, minimum width=9.0cm, minimum height=0.9cm] (e2e) at (\xcenter, \yc) {端到端语音到语音模型};
    \node[font=\small] (sin3) at (\xleft, \yc) {输入语音};
    \node[font=\small] (sout3) at (\xright, \yc) {输出语音};
    \draw[arrow] (sin3) -- (e2e);
    \draw[arrow] (e2e) -- (sout3);
    \node[font=\small\bfseries] at (\xcenter, \yc-0.8) {(c) 端到端模型};
\end{tikzpicture}
\bicaption{\enspace 语音到语音生成的三种建模范式}{\enspace Three modeling paradigms for speech-to-speech generation}
\label{fig:rw_paradigm}
\end{figure}


% ============================================================
% 图 2-2：本章内容结构概览（树状图）
% ============================================================
\begin{figure}[htbp]
\centering
\begin{tikzpicture}[
    node distance=0.6cm and 1.2cm,
    box/.style={draw, rounded corners=3pt, minimum height=0.8cm,
        font=\small, align=center, thick, fill=#1!12, draw=#1!70!black},
    box/.default=blue,
    sec/.style={font=\scriptsize, text=gray!70},
]
    % 根节点 (中心 x=5.5)
    \node[box=violet, minimum width=3.0cm] (root) at (5.5, 0) {语音到语音生成};

    % 第二层：章节号放进框内第二行
    \node[box=teal, minimum width=3.2cm] (encdec) at (2.5, -1.5) {编码器-解码器架构\\[-2pt]{\scriptsize（\S 2.3）}};
    \node[box=orange, minimum width=3.2cm] (deconly) at (8.5, -1.5) {仅解码器架构\\[-2pt]{\scriptsize（\S 2.4）}};

    % 第三层：§编号和任务名都在框内
    % 框宽由文字自动撑开（约3.2cm），四框间距需保证不重叠
    \node[box=teal] (s2tt) at (0.0, -3.8) {语音到文本\\[-2pt]{\scriptsize（\S 2.3.2 语音到文本翻译）}};
    \node[box=blue] (s2st) at (3.7, -3.8) {语音到语音\\[-2pt]{\scriptsize（\S 2.3.3 语音到语音翻译）}};

    \node[box=orange] (s2t_llm) at (7.3, -3.8) {语音到文本\\[-2pt]{\scriptsize（\S 2.4.2 语音理解大模型）}};
    \node[box=red] (s2s_llm) at (11.0, -3.8) {语音到语音\\[-2pt]{\scriptsize（\S 2.4.3 语音交互大模型）}};

    % 连线
    \draw[arrow] (root) -- (encdec);
    \draw[arrow] (root) -- (deconly);
    \draw[arrow] (encdec) -- (s2tt);
    \draw[arrow] (encdec) -- (s2st);
    \draw[arrow] (deconly) -- (s2t_llm);
    \draw[arrow] (deconly) -- (s2s_llm);
\end{tikzpicture}
\bicaption{\enspace 本章内容结构概览}{\enspace Overview of the chapter structure}
\label{fig:rw_overview}
\end{figure}


% ============================================================
% 图 2-5：Transformer 编码器-解码器架构示意图
% ============================================================
\begin{figure}[htbp]
\centering
\begin{tikzpicture}[node distance=0.35cm,
    block/.style={draw, rounded corners=2pt, minimum width=2.4cm, minimum height=0.6cm,
        font=\small, align=center, thick, fill=#1!12, draw=#1!70!black},
    block/.default=blue]

    % 编码器（自底向上）
    \node[block=green, minimum width=2.8cm] (enc_input) at (-2.5, 0) {输入嵌入 + 位置编码};
    \node[block=blue, above=of enc_input, minimum width=2.8cm] (enc_sa) {双向自注意力};
    \node[block=blue, above=of enc_sa, minimum width=2.8cm] (enc_ffn) {前馈网络};
    \draw[arrow2] (enc_input) -- (enc_sa);
    \draw[arrow2] (enc_sa) -- (enc_ffn);
    % N层标注
    \begin{scope}[on background layer]
        \node[dbox=blue!40, fit=(enc_sa)(enc_ffn), inner sep=5pt] (enc_layer) {};
    \end{scope}
    \node[right=0.1cm of enc_layer, font=\scriptsize, blue!70] {$\times N$};

    % 编码器输出
    \node[above=0.35cm of enc_ffn, font=\small] (enc_out) {编码器输出 $Z$};
    \draw[arrow2] (enc_ffn) -- (enc_out);

    % 编码器标题（放在编码器输出上方）
    \node[font=\small\bfseries, above=0.15cm of enc_out] (enc_title) {编码器};

    % 解码器（自底向上）
    \node[block=orange, minimum width=2.8cm] (dec_input) at (3.0, 0) {输出嵌入 + 位置编码};
    \node[block=red, above=of dec_input, minimum width=2.8cm] (dec_sa) {因果自注意力};
    \node[block=violet, above=of dec_sa, minimum width=2.8cm] (dec_ca) {交叉注意力};
    \node[block=red, above=of dec_ca, minimum width=2.8cm] (dec_ffn) {前馈网络};
    \node[block=gray, above=of dec_ffn, minimum width=2.8cm] (dec_linear) {线性层 + Softmax};
    \draw[arrow2] (dec_input) -- (dec_sa);
    \draw[arrow2] (dec_sa) -- (dec_ca);
    \draw[arrow2] (dec_ca) -- (dec_ffn);
    \draw[arrow2] (dec_ffn) -- (dec_linear);
    % N层标注
    \begin{scope}[on background layer]
        \node[dbox=red!40, fit=(dec_sa)(dec_ca)(dec_ffn), inner sep=5pt] (dec_layer) {};
    \end{scope}
    \node[right=0.1cm of dec_layer, font=\scriptsize, red!70] {$\times N$};

    % 输出
    \node[above=0.2cm of dec_linear, font=\small] (output) {输出};
    \draw[arrow2] (dec_linear) -- (output);

    % 解码器标题（放在输出上方）
    \node[font=\small\bfseries, above=0.15cm of output] (dec_title) {解码器};

    % 交叉注意力连接
    \draw[arrow2, violet!70] (enc_out.east) -| ([xshift=-0.3cm]dec_ca.west) -- (dec_ca.west);
\end{tikzpicture}
\bicaption{\enspace Transformer编码器-解码器架构}{\enspace Transformer encoder-decoder architecture}
\label{fig:rw_transformer}
\end{figure}


% ============================================================
% 图 2-6：语音到文本翻译模型结构对比图
% ============================================================
\begin{figure}[htbp]
\centering
\begin{tikzpicture}[node distance=0.3cm and 0.3cm,
    enc/.style={draw, rounded corners=2pt, minimum height=0.7cm, minimum width=1.6cm,
        font=\small, align=center, thick, fill=#1!12, draw=#1!70!black},
    enc/.default=blue,
]
    % 使用绝对坐标，三列左对齐，自底向上构建

    % === (a) 基础编码器-解码器 ===
    \node[enc=green, minimum width=2.4cm] (a_senc) at (0, 0) {语音编码器};
    \node[enc=red, minimum width=2.4cm] (a_dec) at (0, 1.2) {文本解码器};
    \node[font=\scriptsize] at (0, -0.6) {源语言语音};
    \node[font=\scriptsize] at (0, 1.8) {目标语言文本};
    \draw[arrow2] (a_senc) -- (a_dec);
    \node[font=\small\bfseries] at (0, -1.2) {(a) 基础};

    % === (b) 并联双编码器 ===
    \node[enc=green, minimum width=1.5cm] (b_senc) at (3.6, 0) {语音\\编码器};
    \node[enc=orange, minimum width=1.5cm] (b_tenc) at (5.4, 0) {文本\\编码器};
    \node[enc=red, minimum width=2.4cm] (b_dec) at (4.5, 1.8) {文本解码器};
    \node[font=\scriptsize] at (3.6, -0.7) {语音};
    \node[font=\scriptsize] at (5.4, -0.7) {文本};
    \node[font=\scriptsize] at (4.5, 2.4) {目标语言文本};
    \draw[arrow2] (b_senc.north) -- ([xshift=-0.5cm]b_dec.south);
    \draw[arrow2] (b_tenc.north) -- ([xshift=0.5cm]b_dec.south);
    \node[font=\small\bfseries] at (4.5, -1.3) {(b) 并联双编码器};

    % === (c) 串联双编码器 ===
    \node[enc=green, minimum width=2.4cm] (c_senc) at (9.0, 0) {语音编码器};
    \node[enc=orange, minimum width=2.4cm] (c_tenc) at (9.0, 1.0) {文本编码器};
    \node[enc=red, minimum width=2.4cm] (c_dec) at (9.0, 2.0) {文本解码器};
    \node[font=\scriptsize] at (9.0, -0.6) {源语言语音};
    \node[font=\scriptsize] at (9.0, 2.6) {目标语言文本};
    \draw[arrow2] (c_senc) -- (c_tenc);
    \draw[arrow2] (c_tenc) -- (c_dec);
    \node[font=\small\bfseries] at (9.0, -1.2) {(c) 串联双编码器};
\end{tikzpicture}
\bicaption{\enspace 语音到文本翻译的三种模型结构}{\enspace Three model architectures for speech-to-text translation}
\label{fig:rw_s2tt}
\end{figure}


% ============================================================
% 图 2-7：语音到语音翻译流程图
% ============================================================
\begin{figure}[htbp]
\centering
\begin{tikzpicture}[node distance=0.3cm and 0.25cm,
    blk/.style={draw, rounded corners=2pt, minimum height=0.65cm,
        font=\small, align=center, thick, fill=#1!12, draw=#1!70!black},
    blk/.default=blue,
]
    % 使用绝对 y 坐标，每行统一高度；小标题统一在左侧

    % === (a) 一阶段 ===
    \def\ya{0}
    \node[blk=green, minimum width=1.5cm] (a_enc) at (2.5, \ya) {语音\\编码器};
    \node[blk=red, minimum width=1.8cm] (a_dec) at (4.8, \ya) {语音\\解码器};
    \node[blk=teal, minimum width=1.2cm] (a_voc) at (7.2, \ya) {声码器};
    \node[font=\scriptsize] at (1.0, \ya) {源语音};
    \node[font=\scriptsize, anchor=west] at (8.1, \ya) {目标语音};
    \draw[arrow2] (1.4, \ya) -- (a_enc);
    \draw[arrow2] (a_enc) -- (a_dec);
    \draw[arrow2] (a_dec) -- node[above=2pt, slbl]{语音表征} (a_voc);
    \draw[arrow2] (a_voc) -- (7.95, \ya);
    \node[font=\small\bfseries, anchor=east] at (0.3, \ya) {(a) 一阶段};

    % === (b) 两阶段 ===
    \def\yb{-2.0}
    \node[blk=green, minimum width=1.5cm] (b_enc) at (2.5, \yb) {语音\\编码器};
    \node[blk=orange, minimum width=1.5cm] (b_dec1) at (4.3, \yb) {文本\\解码器};
    \node[blk=red, minimum width=1.5cm] (b_dec2) at (6.1, \yb) {语音\\解码器};
    \node[blk=teal, minimum width=1.1cm] (b_voc) at (7.7, \yb) {声码器};
    \node[font=\scriptsize] at (1.0, \yb) {源语音};
    \node[font=\scriptsize, anchor=west] at (8.6, \yb) {目标语音};
    \draw[arrow2] (1.4, \yb) -- (b_enc);
    \draw[arrow2] (b_enc) -- (b_dec1);
    \draw[arrow2] (b_dec1) -- node[above, slbl]{文本} (b_dec2);
    \draw[arrow2] (b_dec2) -- (b_voc);
    \draw[arrow2] (b_voc) -- (8.45, \yb);
    % 标注两个阶段
    \draw[braceup] (b_dec1.north west) -- (b_dec1.north east) node[midway, above=4pt, slbl]{第一阶段};
    \draw[braceup] (b_dec2.north west) -- (b_dec2.north east) node[midway, above=4pt, slbl]{第二阶段};
    \node[font=\small\bfseries, anchor=east] at (0.3, \yb) {(b) 两阶段};

    % === (c) 非自回归 ===
    \def\yc{-4.0}
    \node[blk=green, minimum width=1.5cm] (c_enc) at (2.5, \yc) {语音\\编码器};
    \node[blk=violet, minimum width=1.8cm] (c_dec) at (4.8, \yc) {NAR\\解码器};
    \node[blk=teal, minimum width=1.2cm] (c_voc) at (7.2, \yc) {声码器};
    \node[font=\scriptsize] at (1.0, \yc) {源语音};
    \node[font=\scriptsize, anchor=west] at (8.1, \yc) {目标语音};
    \draw[arrow2] (1.4, \yc) -- (c_enc);
    \draw[arrow2] (c_enc) -- (c_dec);
    \draw[arrow2] (c_dec) -- node[above=2pt, slbl]{语音表征} (c_voc);
    \draw[arrow2] (c_voc) -- (7.95, \yc);
    \node[below=0.15cm of c_dec, font=\scriptsize, violet!70] {并行生成};
    \node[font=\small\bfseries, anchor=east] at (0.3, \yc) {(c) 非自回归};
\end{tikzpicture}
\bicaption{\enspace 语音到语音翻译的三种方法}{\enspace Three approaches for speech-to-speech translation}
\label{fig:rw_s2st}
\end{figure}


% ============================================================
% 图 2-8：GPT 仅解码器架构示意图
% ============================================================
\begin{figure}[htbp]
\centering
\begin{tikzpicture}[node distance=0.35cm,
    block/.style={draw, rounded corners=2pt, minimum width=2.8cm, minimum height=0.6cm,
        font=\small, align=center, thick, fill=#1!12, draw=#1!70!black},
    block/.default=blue]

    \node[block=green, minimum width=3.0cm] (input) at (0, 0) {输入嵌入 + 位置编码};
    \node[block=orange, above=of input, minimum width=3.0cm] (sa) {因果自注意力};
    \node[block=orange, above=of sa, minimum width=3.0cm] (ffn) {前馈网络};
    \node[block=gray, above=0.5cm of ffn, minimum width=3.0cm] (linear) {线性层 + Softmax};

    \draw[arrow2] (input) -- (sa);
    \draw[arrow2] (sa) -- (ffn);
    \draw[arrow2] (ffn) -- (linear);

    % N层标注
    \begin{scope}[on background layer]
        \node[dbox=orange!40, fit=(sa)(ffn), inner sep=5pt] (dec_layer) {};
    \end{scope}
    \node[right=0.1cm of dec_layer, font=\scriptsize, orange!70] {$\times N$};

    % 输入输出
    \node[below=0.2cm of input, font=\small] {输入序列};
    \node[above=0.2cm of linear, font=\small] {下一词元预测};

\end{tikzpicture}
\bicaption{\enspace 仅解码器架构}{\enspace Decoder-only architecture}
\label{fig:rw_decoderonly}
\end{figure}


% ============================================================
% 图 2-9：语音到文本生成架构图（编码器 + 适配器 + LLM）
% ============================================================
\begin{figure}[htbp]
\centering
\begin{tikzpicture}[node distance=0.3cm and 0.4cm,
    blk/.style={draw, rounded corners=2pt, minimum height=0.8cm,
        font=\small, align=center, thick, fill=#1!12, draw=#1!70!black},
    blk/.default=blue,
]
    \node[blk=green, minimum width=2.2cm] (senc) {语音编码器\\{\scriptsize (Whisper 等)}};
    \node[blk=orange, right=1.3cm of senc, minimum width=2.0cm] (adapt) {适配器\\{\scriptsize (MLP / Q-Former / CTC)}};
    \node[blk=red, right=0.8cm of adapt, minimum width=2.5cm] (llm) {大语言模型\\{\scriptsize (LLaMA / Qwen 等)}};

    \node[left=0.3cm of senc, font=\small] (in) {语音};
    \node[right=0.3cm of llm, font=\small] (out) {文本};

    \draw[arrow] (in) -- (senc);
    \draw[arrow] (senc) -- node[above, slbl]{语音表示} (adapt);
    \draw[arrow] (adapt) -- node[above, slbl]{嵌入} (llm);
    \draw[arrow] (llm) -- (out);

    % 标注冻结/可训练
    \node[below=0.2cm of senc, font=\scriptsize, green!60!black] {通常冻结};
    \node[below=0.2cm of adapt, font=\scriptsize, orange!60!black] {可训练};
    \node[below=0.2cm of llm, font=\scriptsize, red!60!black] {微调 / LoRA};
\end{tikzpicture}
\bicaption{\enspace 基于大语言模型的语音到文本生成架构}{\enspace Architecture of LLM-based speech-to-text generation}
\label{fig:rw_s2t_llm}
\end{figure}


% ============================================================
% 图 2-10：语音到语音生成架构图（原生式 vs 组合式）
% ============================================================
\begin{figure}[htbp]
\centering
\begin{tikzpicture}[node distance=0.3cm and 0.35cm,
    blk/.style={draw, rounded corners=2pt, minimum height=0.7cm,
        font=\small, align=center, thick, fill=#1!12, draw=#1!70!black},
    blk/.default=blue,
]
    % === (a) 原生式 ===
    \def\ya{0}
    \node[blk=violet, minimum width=5.0cm, minimum height=1.2cm] (native) at (4.0, \ya) {
        语音-语言模型\\
        {\scriptsize 统一词表：文本 token + 语音 token}
    };
    \node[left=0.3cm of native, font=\small] (in1) {语音};
    \node[right=0.3cm of native, font=\small] (out1) {语音};
    \draw[arrow] (in1) -- (native);
    \draw[arrow] (native) -- (out1);
    \node[font=\small\bfseries] at (4.0, \ya+1.2) {(a) 原生式};

    % === (b) 组合式 ===
    \def\yb{-2.5}
    \node[blk=green, minimum width=1.6cm] (b_senc) at (1.0, \yb) {语音\\编码器};
    \node[blk=orange, minimum width=1.2cm] (b_adapt) at (3.0, \yb) {适配器};
    \node[blk=red, minimum width=1.8cm] (b_llm) at (4.8, \yb) {大语言模型};
    \node[blk=blue, minimum width=1.6cm] (b_sdec) at (6.8, \yb) {语音\\解码器};

    \node[left=0.3cm of b_senc, font=\small] (in2) {语音};
    \node[right=0.3cm of b_sdec, font=\small] (out2) {语音};

    \draw[arrow2] (in2) -- (b_senc);
    \draw[arrow2] (b_senc) -- (b_adapt);
    \draw[arrow2] (b_adapt) -- (b_llm);
    \draw[arrow2] (b_llm) -- (b_sdec);
    \draw[arrow2] (b_sdec) -- (out2);

    % 文本输出
    \node[above=0.3cm of b_llm, font=\scriptsize] (txt_out) {文本};
    \draw[arrow2] (b_llm) -- (txt_out);

    \node[below=0.15cm of b_sdec, font=\scriptsize, blue!60] {NAR / AR};
    \node[font=\small\bfseries] at (4.0, \yb+1.5) {(b) 组合式};
\end{tikzpicture}
\bicaption{\enspace 语音到语音生成的两种架构：原生式与组合式}{\enspace Two architectures for speech-to-speech generation: native vs.\ modular}
\label{fig:rw_s2s_llm}
\end{figure}
